\section{Model formulation}
\label{sec:model}

\subsection{Governing equations}
\label{sec:equations}

The water has constant density $\rho$ and occupies
$z_b(x,y,t) \le z \le \eta(x,y,t)$, where $z_b$ is the seabed and $\eta$ the
free surface. With the velocity $\bm{u}=(u,v,w)$, the incompressible
Navier--Stokes equations read
\begin{align}
  \nabla\cdot\bm{u} &= 0, \label{eq:cont}\\
  \frac{\partial \bm{u}}{\partial t} + \nabla\cdot(\bm{u}\bm{u})
  - \nabla\cdot\left[\nu_{\mathrm{eff}}\left(\nabla\bm{u} + \nabla\bm{u}^{T}\right)\right]
  &= -\nabla \prgh + \bm{f}, \label{eq:mom}
\end{align}
where $\nu_{\mathrm{eff}}$ is the molecular plus any turbulent viscosity and
$\bm{f}$ the forcing of the relaxation zones (Section~\ref{sec:zones}). The
kinematic pressure is written as
\begin{equation}
  \prgh = \frac{p}{\rho} + g z = g\eta + q ,
  \label{eq:prgh}
\end{equation}
where $g\eta$ is the hydrostatic part and $q$ the non-hydrostatic part. In
hydrostatic equilibrium $p/\rho = g(\eta - z)$, so $\prgh = g\eta$ is uniform in
the vertical. Gravity does not appear explicitly in \eqref{eq:mom}, so still
water stays exactly at rest on any mesh, however skewed the cells are over a
steep bottom.

\subsection{$\sigma$-mesh and its motion}
\label{sec:sigma}

The mesh is a vertical extrusion of a horizontal mesh. Each face of the top
patch (\code{surface}) defines one water column, and the points of a column lie
on vertical lines. The points move vertically only:
\begin{equation}
  z = z_b + \sigma\,(\eta - z_b), \qquad \sigma\in[0,1],
  \label{eq:sigma}
\end{equation}
where $\sigma$ of each point is taken from the initial mesh, so any vertical
grading is kept. The equations are solved in physical coordinates on the moving
mesh (arbitrary Lagrangian--Eulerian form): there are no metric terms of a
$\sigma$-transformation, and OpenFOAM's swept-volume mesh fluxes satisfy the
space conservation law. The elevation $\eta$ is stored once per column and
interpolated to the surface points with inverse-distance weights.

\subsection{Free surface and bottom}
\label{sec:bc}

At the free surface the pressure is atmospheric, so
\begin{equation}
  q = 0, \qquad \prgh = g\eta \qquad \text{at } z=\eta .
  \label{eq:dynamic}
\end{equation}
The kinematic condition is used in its integral form: the change of the
water depth of a column balances the net inflow,
\begin{equation}
  \frac{\partial (\eta - z_b)}{\partial t} + \nabla_h\cdot \int_{z_b}^{\eta}\bm{u}_h\,\dd z = 0 ,
  \label{eq:kinematic}
\end{equation}
which the discrete method imposes column by column
(Section~\ref{sec:freesurface}). The bottom is impermeable, with a slip or
no-slip condition. When the bottom moves, the boundary condition
\code{nhMovingSlipVelocity} gives the fluid the normal velocity of the bottom
and leaves it free to slip tangentially, so that the flux through the bottom
faces equals their swept volume. A quadratic bottom stress,
$\bm{\tau}_b/\rho = C_d\,|\bm{u}_b|\,\bm{u}_b$ with $\bm{u}_b$ the velocity of
the bottom cell and $C_d$ constant, $g n^2/D^{1/3}$ from a Manning coefficient,
or $[\kappa/\ln(30z_1/k_s)]^2$ from the log law at the centre $z_1$ of the
bottom cell with roughness height $k_s$ (NHWAVE's \code{Zob}), can be added.
Lateral boundaries are walls, symmetry planes, the outer edges of relaxation
zones or periodic (cyclic) pairs. Moving shorelines (wetting and drying) are described in
Section~\ref{sec:wd}.

\subsection{Wave generation and absorption}
\label{sec:zones}

Relaxation zones relax the velocity and the surface towards a target
$(\bm{u}_t,\eta_t)$:
\begin{equation}
  \bm{f} = -\sigma_r(\chi)\,(\bm{u}-\bm{u}_t),\qquad
  \frac{\partial \eta}{\partial t} = \ldots - \sigma_r(\chi)\,(\eta-\eta_t),\qquad
  \sigma_r(\chi) = \sigma_{\max}\,\chi^{n},
  \label{eq:relax}
\end{equation}
where $\chi\in[0,1]$ grows linearly from the inner edge of the zone
(\code{start}) to its outer edge (\code{end}). Both terms are implicit. A zone of
type \code{wave} uses a regular wave of height $H$ and period $T$ as the
target, ramped up over a time $t_{\mathrm{ramp}}$; it generates the incident
waves and absorbs the reflected ones. The wave is linear (Airy), or a
stream-function wave of any height \citep{fenton1988}, whose Fourier
coefficients are computed beforehand (Section~\ref{sec:wd-method}). A zone of type \code{sponge} has still
water as its target, or relaxes the velocity only (\code{relaxEta false}),
which keeps the water volume, e.g.\ ahead of a truncated beach.

\subsection{Seabed motion}
\label{sec:bottommotion}

A prescribed displacement moves the bottom points,
$z_b(x,y,t) = z_{b0}(x,y) + \Delta z(x,y,t)$, and the $\sigma$-mesh follows
\eqref{eq:sigma}. No change of the free-surface condition is needed, since the
surface flux of each column already contains the volume pushed up by the
bottom. Two motions are built in.
\begin{itemize}[nosep]
  \item \code{uplift}: a box raised by $\zeta_0$ with an exponential,
        $\zeta_0\,(1-e^{-\alpha t})$ with $\alpha = 1.11/t_c$
        \citep{hammack1973}, a half-cosine or a linear history.
  \item \code{slide}: a rigid quasi-Gaussian slide on a plane slope
        \citep{grilli2005,enet2007}, of thickness
        $\zeta = \frac{T_s}{1-\epsilon}\left[\mathrm{sech}(k_b\xi)\,\mathrm{sech}(k_w\chi)-\epsilon\right]$
        normal to the slope, with $k_b = 2C/b$, $k_w = 2C/w$ and
        $C = \mathrm{acosh}(1/\epsilon)$, moving down the slope as
        $s(t) = s_0\ln\cosh(t/t_0)$ with $s_0 = u_t^2/a_0$ and
        $t_0 = u_t/a_0$, where $u_t$ is the terminal velocity and $a_0$ the
        initial acceleration.
\end{itemize}

\subsection{Breaking}
\label{sec:hfa}

A column starts breaking when its surface rises fast,
\begin{equation}
\frac{\partial\eta}{\partial t} > \alpha\sqrt{gD},
\label{eq:onset}
\end{equation}
with $D$ the local depth, or, optionally, when the horizontal velocity of its
surface cell exceeds a fraction of the crest speed, estimated by $\sqrt{gD}$:
$|\bm{u}_s| > B_k\sqrt{gD}$ (kinematic onset; $B_k = 0.85$, the ratio $U/C$
at the onset of breaking found by \citet{barthelemy2018}). A breaking column
keeps breaking, and a neighbouring column joins it, while
\begin{equation}
\frac{\partial\eta}{\partial t} > \eta_t^*,\qquad
\eta_t^* = \left[\beta + (\beta_F - \beta)\min\!\left(\frac{t-t_0}{T^*},1\right)\right]\sqrt{gD},
\qquad T^* = T_s\sqrt{D/g}.
\label{eq:persist}
\end{equation}
Without $\beta_F$ the threshold is the constant $\beta\sqrt{gD}$ of
\citet{smit2013} ($\alpha = 0.6$, $\beta = 0.3$). With $\beta_F$ it falls with
the age $t-t_0$ of the breaking event, as in the eddy-viscosity model of
\citet{kennedy2000} ($\beta_F = 0.15$, $T_s = 5$; Kennedy et al.\ start at the
onset threshold, $\beta = \alpha$). A column that joins next to an older
breaker takes over its age, so the age travels with the wave and the
breaking region grows from the front over the whole face of a mature bore.
Neighbours include the columns across cyclic patches and processor
boundaries. Two treatments act in
the breaking columns, alone or together (\code{model HFA},
\code{eddyViscosity} or \code{hybrid}).
\begin{itemize}[nosep]
  \item \textbf{Hydrostatic front approximation} \citep{smit2013}. The
        non-hydrostatic pressure is switched off, so the wave front travels as
        a hydrostatic bore and the conservative momentum equation with a
        shock-capturing (TVD) convection scheme dissipates its energy. How $q$
        is switched off is described in Section~\ref{sec:numhfa}.
  \item \textbf{Breaking eddy viscosity} \citep{kennedy2000}. The momentum
        equation gets the diffusion term $\nabla\cdot(\nu_b\nabla\bm{u})$ with
        \begin{equation}
        \nu_b = B\,\delta_b^2\,D\,\frac{\partial\eta}{\partial t},\qquad
        B = \min\!\left[\max\!\left(\frac{\partial\eta/\partial t}{\eta_t^*}-1,\,0\right),1\right],
        \label{eq:nub}
        \end{equation}
        $\delta_b = 1.2$, uniform over the depth of the column. It acts through
        the side faces of the columns (horizontal mixing), optionally also
        through the layer faces (\code{verticalMixing}), and not through the
        closed faces of dry columns. The non-hydrostatic pressure is kept.
\end{itemize}
Subgrid turbulence can be added with OpenFOAM's turbulence library. The
Smagorinsky model of the LES library with the filter width $V^{1/3}$ of the
cell is the model NHWAVE uses for the longshore current \citep{ma2012},
$\nu_t = (C_s\Delta)^2\sqrt{2S_{ij}S_{ij}}$ with $\Delta = (\Delta x\,\Delta y\,
\Delta\sigma D)^{1/3}$; its default constants give $C_s = 0.168$.
